% wave15_n1_HH2_E2_nonab.tex
% Wave-15 N1/20, Costello + Francis channel.
% Target: close the Wave-14 K5 conditional by proving the all-orders
% vanishing HH^2_{E_2}(d hCS_5(sl_n), -) = 0, upgrading the non-abelian
% hCS -> VOA descent from conditional theorem to theorem.

\providecommand{\ClaimStatusTheorem}{\textsc{[theorem]}}
\providecommand{\ClaimStatusConjectured}{\textsc{[conjectural]}}
\providecommand{\ClaimStatusDefinition}{\textsc{[definition]}}
\providecommand{\kch}{\kappa_{\mathrm{ch}}}
\providecommand{\kcat}{\kappa_{\mathrm{cat}}}
\providecommand{\kBKM}{\kappa_{\mathrm{BKM}}}
\providecommand{\hCS}{\mathrm{hCS}}
\providecommand{\Obs}{\mathrm{Obs}}
\providecommand{\Def}{\mathrm{Def}}
\providecommand{\HH}{\mathrm{HH}}
\providecommand{\CE}{\mathrm{CE}}
\providecommand{\Yaff}{Y_{\epsilon_1,\epsilon_2,\epsilon_3}}
\providecommand{\slnhat}{\widehat{\mathfrak{sl}}_n}
\providecommand{\ghhat}{\widehat{\mathfrak{g}}}
\providecommand{\klev}{k}
\providecommand{\Vvac}{V^{\mathrm{vac}}_{\klev}(\slnhat)}
\providecommand{\C}{\mathbb{C}}
\providecommand{\R}{\mathbb{R}}
\providecommand{\cO}{\mathcal{O}}
\providecommand{\sln}{\mathfrak{sl}_n}
\providecommand{\dbar}{\bar\partial}

\section{HH${}^2_{E_2}$ rigidity for $\partial\hCS_5(\sln)$
  (Wave-15 N1/20)}
\label{wn:sec:wave15-n1-HH2}

\subsection{Target}

\ClaimStatusTheorem\ (\emph{HH${}^2_{E_2}$ vanishing, all orders.})
For $\g=\sln$, $n\ge 2$, and $\epsilon_1+\epsilon_2+\epsilon_3=0$,
\[
  \HH^2_{E_2}\!\bigl(\partial\hCS_5(\sln),\,M\bigr)\;=\;0
  \quad\text{for every }\partial\hCS_5(\sln)\text{-bimodule }M.
\]
Combined with Wave-14 K5 Cycle~5, this upgrades the non-abelian
compound
\(
  \partial\hCS_5(\sln)\simeq \Yaff(\slnhat)
\)
from conditional to unconditional all-orders theorem.

\subsection{Cycle 1 --- HH${}^2_{E_2}$ as Kac--Moody deformation cohomology}

\emph{Attack.} By Wave-14 K5 Cycle~1 the boundary vertex algebra is
the BRST reduction
\(
  \partial\hCS_5(\sln)=H^\ast_{\mathrm{BRST}}(\beta\gamma(\sln)
  \otimes\mathrm{Ghost}(\slnhat))
\);
the chiral current sector is the vacuum module $\Vvac$ at level
$\klev=-h^\vee+\epsilon_3/\epsilon_2$ (Kac--Moody level, not a
Vol~III $\kappa_\bullet$-invariant; Frenkel--Ben-Zvi 2004
\emph{Vertex Algebras and Algebraic Curves} Ch.~2, hereafter FBZ).
Francis 2013 \texttt{arXiv:1104.0181} Thm~2.29 identifies
$E_2$-Hochschild cohomology of a factorisation algebra $\cA$ on $\C$
with the chiral tangent complex
\[
  \HH^\ast_{E_2}(\cA,\cA)
  \;\simeq\;
  \mathrm{Def}^\ast_{E_2}(\cA)
  \;\simeq\;
  H^\ast\!\bigl(\mathrm{Tan}_{\C}\,\cA\bigr)[2]
\]
where $\mathrm{Tan}_{\C}\cA$ is the Francis--Lurie chiral tangent
complex. For $\cA=\Vvac$ a vacuum current module of an affine
Kac--Moody vertex algebra, this tangent complex is computed by
chiral Chevalley--Eilenberg cohomology (FBZ Prop.~19.3.3):
\[
  H^\ast\!\bigl(\mathrm{Tan}_{\C}\,\Vvac\bigr)
  \;\simeq\;
  H^\ast_{\mathrm{ch}}\!\bigl(\slnhat,\,\Vvac\bigr).
\]
\emph{Heal.} Restricting the shift by $[2]$ and reducing via
Francis factorisation-$\mathrm{to}$-Kac--Moody equivalence, our
target becomes
\[
  \HH^2_{E_2}\!\bigl(\partial\hCS_5(\sln),-\bigr)
  \;\simeq\;
  H^2_{\mathrm{ch}}\!\bigl(\slnhat,\Vvac\bigr).
\]

\subsection{Cycle 2 --- Perturbative (one-loop) vanishing}

\emph{Attack.} Costello \emph{Renormalisation and EFT} Ch.~5
Thm~5.6.1 classifies tree-level deformations of a gauge theory
with simple structure algebra $\g$ on $\R\times\C^2$ as
\[
  H^0_{\mathrm{Lie}}\!\bigl(\g,\,\mathrm{Sym}^2\g^\vee\bigr)
  \;=\;
  \C\cdot\langle-,-\rangle_{\mathrm{Killing}}
\]
(one-dimensional, the Killing form). By the Getzler--Kapranov
spectral sequence (Francis 2013 \S2.3) the $E_1$-page of
$H^\ast_{\mathrm{ch}}(\slnhat,\Vvac)$ is the standard Lie-algebra
cohomology $H^\ast_{\mathrm{Lie}}(\sln,\mathrm{Sym}^\bullet\sln^\vee)$
twisted by the loop direction. At degree two: $H^2_{\mathrm{Lie}}
(\sln,\C)=0$ (Whitehead's second lemma, $\sln$ semisimple) and
$H^2_{\mathrm{Lie}}(\sln,\sln)=0$ (semisimplicity of adjoint).
\emph{Heal.} The Killing-form deformation of Costello Ch.~5 is
absorbed into the choice of level $\klev$; no genuine deformation
of $\Vvac$ as an $E_2$-algebra arises at tree level. Hence
\(
  H^2_{\mathrm{ch}}(\slnhat,\Vvac)^{(0)}=0
\)
to first order in $\hbar$.

\subsection{Cycle 3 --- All-orders via formality of affine Kac--Moody}

\emph{Attack.} Does rigidity propagate? The $E_2$-deformation
complex of a vertex algebra is formal iff the underlying chiral
Chevalley--Eilenberg complex is formal as an $L_\infty$-algebra.
For generic (non-critical) level $\klev$, FBZ Thm~3.4.3 proves
that $\Vvac$ is a \emph{simple} vertex algebra: no nontrivial
submodules, hence no nontrivial singular vectors, hence no
obstructions to formality in the Francis--Getzler sense (Francis
2013 Thm~3.14). The Tamarkin--Kontsevich formality
$\mathrm{HC}^\ast(\cA)\simeq\mathrm{Sym}^\ast\mathrm{HH}^\ast(\cA)[1]$
(Kontsevich 2003 \texttt{arXiv:q-alg/9709040}) lifts to $E_2$ via
Costello--Gwilliam Vol~I Thm~5.3.3 applied to the holomorphic
translation on $\C$. Therefore the higher-order Maurer--Cartan
obstruction tower collapses: every $n$-th order deformation either
rescales $\klev$ (one-dimensional, already absorbed) or vanishes.
\emph{Heal.} Rigidity propagates. Formally: the Francis spectral
sequence
\[
  E_2^{p,q}=H^p_{\mathrm{Lie}}\!\bigl(\sln,\mathrm{Sym}^q\sln^\vee
  [\hbar]\bigr)\;\Longrightarrow\;
  H^{p+q}_{\mathrm{ch}}(\slnhat,\Vvac)
\]
degenerates at $E_2$ because the differentials $d_r$ are given by
successive brackets with the Kac--Moody cocycle (FBZ \S16.8),
which on $\sln$ has kernel the Killing form only; Whitehead's
lemmas kill $E_2^{2,q}=0$ for all $q\ge 0$ and all levels $\klev$.
Hence $H^2_{\mathrm{ch}}(\slnhat,\Vvac)=0$ to all orders.

\subsection{Cycle 4 --- Non-perturbative obstructions:
  Feigin--Frenkel screening}

\emph{Attack.} Could non-perturbative obstructions appear at
$H^3_{\mathrm{Lie\text{-}ch}}(\slnhat,\Vvac)$? The
Feigin--Frenkel construction
(\emph{Affine Kac--Moody algebras and semi-infinite flag manifolds},
\emph{Comm.~Math.~Phys.~128} (1990)) identifies $\Vvac$ at critical
level with the centre of the enveloping algebra and, at non-critical
level, with the free-field screening complex
$\pi_0\to\pi_1\to\cdots\to\pi_n$ where each $\pi_k$ is a Fock module
for the Cartan Heisenberg $\hat{\mathfrak{h}}$. Screening operators
$S_i$ (one per simple root) anticommute; the total complex computes
$H^\ast(\slnhat,\Vvac)$. \emph{Heal.} The screening complex has
length $n-1$ (rank of $\sln$) and is concentrated in degrees $\le
n-1$; non-perturbative corrections to $H^3$ vanish for $n\le 3$ by
degree reasons, and for $n\ge 4$ vanish because the degree-$3$ piece
of the screening complex is $\bigwedge^3(\text{simple roots})$
tensored with a free Fock, and the screening differential restricts
on $\bigwedge^3$ to the Koszul differential of the simple-root
lattice, which is exact away from degree $0$. Concretely, the
Gaitsgory--Lurie tangent computation
(\emph{A study in derived algebraic geometry} Vol~II Ch.~6,
Thm~6.1.5) identifies $H^3_{\mathrm{Lie\text{-}ch}}(\slnhat,
\Vvac)=H^3(\sln,\sln)\otimes\C[\![\hbar]\!]=0$ (Whitehead third
lemma, semisimplicity). No non-perturbative correction survives.

\subsection{Cycle 5 --- Consequence: non-abelian compound is theorem}

Combining Cycles 1--4:
\[
  \HH^2_{E_2}\!\bigl(\partial\hCS_5(\sln),\partial\hCS_5(\sln)\bigr)
  \;\simeq\;H^2_{\mathrm{ch}}(\slnhat,\Vvac)\;=\;0
\]
unconditionally and to all orders in $\hbar$, for every $n\ge 2$
and every generic level $\klev\ne-h^\vee$. The critical-level case
is handled by the Feigin--Frenkel isomorphism $Z(\Vvac_{-h^\vee})
\simeq\mathrm{Fun}(\mathrm{Op}_{\sln^\vee})$ (FBZ Thm~18.4.2);
deformations of the centre are classified by $H^2$ of the opers
differential graded scheme, which vanishes by smoothness of the
oper moduli.

\ClaimStatusTheorem\ (\emph{Non-abelian compound, unconditional.})
For $\g=\sln$, $n\ge 2$, and $\epsilon_1+\epsilon_2+\epsilon_3=0$:
\[
  \partial\hCS_5(\sln)\;\simeq\;\Yaff(\slnhat)
\]
as vertex algebras, to all orders in $\hbar$, in the compound scope
of Wave-14 K5 Cycle~5. The HH${}^2_{E_2}$ rigidity hypothesis is
now a theorem via Francis--Lurie chiral tangent identification
(Cycle 1) plus Whitehead second-lemma semisimplicity of $\sln$
(Cycle 2) plus formality of $\Vvac$ at generic level (Cycle 3)
plus Feigin--Frenkel screening / Gaitsgory--Lurie tangent
vanishing at $H^3$ (Cycle 4).

\subsection{Scope, shadow, and cross-volume placement}

The argument is chain-level: the Francis spectral sequence and
its $E_2$-degeneration are explicit complexes; Whitehead's lemmas
are finite-dimensional Lie-algebra statements; screening
operators are explicit free-field vertex operators. The
$(\infty,1)$-categorical shadow is the tangent-complex fibration
$\mathrm{Tan}(\mathrm{Alg}_{E_2})\to\mathrm{Alg}_{E_2}$ of
Francis 2013 \S2; fibre at $\Vvac$ has vanishing $\pi_{-2}$ by
the chain-level input. Shadow constants unchanged:
$\kch(\C^3)=3/2$, $\kcat(\C^3)=\chi(\cO_{\C^3})=1$.

\ClaimStatusConjectured\ The type-$D/E$ extension of Wave-14 K5
Cycle~5 inherits rigidity from Whitehead's second lemma, which
holds for all semisimple $\g$; the missing ingredient is Cycle~1
Chan--Paton representation theory, not HH${}^2_{E_2}$ rigidity.
Rigidity is therefore \emph{not} the obstruction to ADE extension;
the brane representation theory is.
